% Topic 1.1.09 — Differentiation. Taylor's Formula.

\section{Differentiation. Taylor's Formula}
\label{sec:1.1.09-derivatives-taylor}

\begin{ticket}
Differentiation of functions. Application of the derivative to finding extrema. Taylor's formula.
\end{ticket}

\subsection*{Theory}

We work with real-valued functions of a single real variable defined
on an open subset of~$\R$. Throughout, $x_0$ denotes a chosen interior
point of the domain.

\begin{definition}
\label{def:derivative}
A function $f$ is \emph{differentiable at~$x_0$} if the limit
\[
  f'(x_0) \;=\; \lim_{h \to 0} \frac{f(x_0 + h) - f(x_0)}{h}
\]
exists and is finite. Equivalently, there exists $A \in \R$ for which
\[
  f(x_0 + h) \;=\; f(x_0) + A \cdot h + o(h)
  \qquad \text{as } h \to 0,
\]
in which case $A = f'(x_0)$. The function~$f$ is \emph{differentiable
on an open set $U$} if it is differentiable at every $x_0 \in U$, and
the resulting function $x \mapsto f'(x)$ is the \emph{derivative}.
\end{definition}

The standard differentiation rules — linearity, product, quotient,
chain — are taken as known and follow directly from
\cref{def:derivative} by manipulating difference quotients. We focus
on what the derivative says about the function.

\begin{proposition}[Differentiability implies continuity]
\label{prop:diff-implies-cont}
If~$f$ is differentiable at~$x_0$, then $f$ is continuous at~$x_0$.
\end{proposition}
\proofof{diff-implies-cont}

\begin{definition}
\label{def:local-extremum}
A function~$f$ has a \emph{local maximum} (resp.\ \emph{local minimum})
at~$x_0$ if there is a neighbourhood $U \ni x_0$ on which
$f(x) \le f(x_0)$ (resp.\ $f(x) \ge f(x_0)$) for all $x \in U$. A
\emph{local extremum} is either a local maximum or a local minimum.
The extremum is \emph{strict} if equality only holds at $x = x_0$.
\end{definition}

\begin{theorem}[Fermat]
\label{thm:fermat}
If~$f$ has a local extremum at an interior point~$x_0$ and is
differentiable at~$x_0$, then $f'(x_0) = 0$.
\end{theorem}
\proofof{fermat}

\begin{theorem}[Rolle]
\label{thm:rolle}
Let~$f$ be continuous on $[a, b]$ and differentiable on $(a, b)$,
with $f(a) = f(b)$. Then there exists $c \in (a, b)$ with
$f'(c) = 0$.
\end{theorem}
\proofof{rolle}

\begin{theorem}[Lagrange's mean-value theorem]
\label{thm:lagrange-mvt}
Let~$f$ be continuous on $[a, b]$ and differentiable on $(a, b)$.
There exists $c \in (a, b)$ with
\[
  f(b) - f(a) \;=\; f'(c)\,(b - a).
\]
\end{theorem}
\proofof{lagrange-mvt}

\begin{theorem}[First-derivative test for extrema]
\label{thm:first-deriv-test}
Let~$f$ be continuous on a neighbourhood of~$x_0$ and differentiable
on a punctured neighbourhood. If $f'(x) > 0$ for $x < x_0$ near~$x_0$
and $f'(x) < 0$ for $x > x_0$ near~$x_0$, then~$f$ has a strict
local maximum at~$x_0$. The mirror sign pattern gives a strict local
minimum.
\end{theorem}
\proofof{first-deriv-test}

\begin{theorem}[Taylor's formula with Peano remainder]
\label{thm:taylor-peano}
Let~$f$ be $n$ times differentiable at~$x_0$. Then as $x \to x_0$,
\[
  f(x) \;=\; \sum_{k=0}^{n} \frac{f^{(k)}(x_0)}{k!}\,(x - x_0)^k \;+\; o\bigl((x - x_0)^n\bigr).
\]
\end{theorem}
\proofof{taylor-peano}

\begin{theorem}[Taylor's formula with Lagrange remainder]
\label{thm:taylor-lagrange}
Let~$f$ be $(n+1)$ times differentiable on an open interval
containing $x_0$ and~$x$. There exists~$c$ strictly between $x_0$
and~$x$ such that
\[
  f(x) \;=\; \sum_{k=0}^{n} \frac{f^{(k)}(x_0)}{k!}\,(x - x_0)^k
  \;+\; \frac{f^{(n+1)}(c)}{(n+1)!}\,(x - x_0)^{n+1}.
\]
\end{theorem}
\proofof{taylor-lagrange}

\begin{corollary}[Second-derivative test]
\label{cor:second-deriv-test}
Suppose~$f$ is twice differentiable at~$x_0$ with $f'(x_0) = 0$.
\begin{itemize}
  \item If $f''(x_0) > 0$, then~$f$ has a strict local minimum at~$x_0$.
  \item If $f''(x_0) < 0$, then~$f$ has a strict local maximum at~$x_0$.
  \item If $f''(x_0) = 0$, the test is inconclusive.
\end{itemize}
\end{corollary}
\proofof{second-deriv-test}

\begin{remark}
The proofs below use only \cref{thm:rolle,thm:lagrange-mvt} and
the Cauchy generalized MVT (essentially Rolle applied to a
two-component auxiliary function). Comprehensive treatment, including
Taylor with the integral form of the remainder and consequences for
real-analytic functions, is in \cite{Zorich1981}.
\end{remark}

\subsection*{Proofs}

\begin{proof}[\Cref{prop:diff-implies-cont}]
\proofanchor{diff-implies-cont}
For $h \neq 0$,
\[
  f(x_0 + h) - f(x_0) \;=\; \frac{f(x_0 + h) - f(x_0)}{h} \cdot h.
\]
The right-hand side has limit $f'(x_0) \cdot 0 = 0$ as $h \to 0$, so
$f(x_0 + h) \to f(x_0)$, which is continuity at~$x_0$.
\end{proof}

\begin{proof}[\Cref{thm:fermat}]
\proofanchor{fermat}
Suppose~$x_0$ is a local maximum (the minimum case is mirrored). For
all~$h$ small enough, $f(x_0 + h) \le f(x_0)$, so
$f(x_0 + h) - f(x_0) \le 0$. Therefore
\[
  \frac{f(x_0 + h) - f(x_0)}{h} \;\le\; 0 \text{ for } h > 0,
  \qquad
  \frac{f(x_0 + h) - f(x_0)}{h} \;\ge\; 0 \text{ for } h < 0.
\]
Both one-sided limits exist and equal $f'(x_0)$ by the assumed
differentiability, so $f'(x_0) \le 0$ and $f'(x_0) \ge 0$
simultaneously, giving $f'(x_0) = 0$.
\end{proof}

\begin{proof}[\Cref{thm:rolle}]
\proofanchor{rolle}
Continuity on the compact interval $[a, b]$ gives $f$ a maximum
and a minimum, attained at some $c_M, c_m \in [a, b]$. If both lie
at the endpoints, $f(a) = f(b)$ forces the maximum to equal the
minimum, so~$f$ is constant; then $f'(c) = 0$ for any $c \in (a, b)$.
Otherwise at least one of $c_M, c_m$ is interior to $(a, b)$, and
applying \cref{thm:fermat} at that point yields $f'(c) = 0$.
\end{proof}

\begin{proof}[\Cref{thm:lagrange-mvt}]
\proofanchor{lagrange-mvt}
Apply \cref{thm:rolle} to the auxiliary function
\[
  g(x) \;=\; f(x) - f(a) - \frac{f(b) - f(a)}{b - a}\,(x - a),
\]
which satisfies $g(a) = 0 = g(b)$ and is continuous on $[a, b]$ /
differentiable on $(a, b)$. Rolle gives some $c \in (a, b)$ with
$g'(c) = 0$, i.e.\ $f'(c) = (f(b) - f(a))/(b - a)$.
\end{proof}

\begin{proof}[\Cref{thm:first-deriv-test}]
\proofanchor{first-deriv-test}
Consider the case $f' > 0$ on the left of~$x_0$ and $f' < 0$ on the
right (the mirror sign pattern is handled symmetrically). Choose
$\delta > 0$ such that~$f$ is continuous on $(x_0 - \delta, x_0 + \delta)$,
differentiable on $(x_0 - \delta, x_0) \cup (x_0, x_0 + \delta)$, with
$f' > 0$ on the left and $f' < 0$ on the right.

For $x \in (x_0 - \delta, x_0)$, apply the mean-value theorem
(\cref{thm:lagrange-mvt}) to~$f$ on $[x, x_0]$: there is
$\xi \in (x, x_0)$ with $f(x_0) - f(x) = f'(\xi)(x_0 - x)$. Both
$f'(\xi) > 0$ (since $\xi$ is on the left) and $x_0 - x > 0$, so
$f(x_0) > f(x)$.

For $x \in (x_0, x_0 + \delta)$, applying the mean-value theorem on
$[x_0, x]$ gives $\eta \in (x_0, x)$ with $f(x) - f(x_0) = f'(\eta)(x - x_0)$.
Now $f'(\eta) < 0$ and $x - x_0 > 0$, so $f(x) < f(x_0)$.

In both cases $f(x_0) > f(x)$ for $x \neq x_0$ with $|x - x_0| < \delta$,
witnessing a strict local maximum.
\end{proof}

\begin{proof}[\Cref{thm:taylor-peano}]
\proofanchor{taylor-peano}
Let $T_n(x) = \sum_{k=0}^{n} \frac{f^{(k)}(x_0)}{k!} (x - x_0)^k$ be
the Taylor polynomial and $r_n(x) = f(x) - T_n(x)$. We must show
$r_n(x) = o((x - x_0)^n)$ as $x \to x_0$.

By construction $r_n(x_0) = 0$ and $r_n^{(k)}(x_0) = 0$ for
$k = 0, 1, \dots, n$ (the $k$-th derivatives of $T_n$ and~$f$
agree at $x_0$ up to order~$n$).

Apply Cauchy's generalized MVT iteratively. For each~$x$ near~$x_0$
there exists~$\xi_1$ strictly between $x_0$ and~$x$ with
\[
  \frac{r_n(x)}{(x - x_0)^n}
  \;=\; \frac{r_n(x) - r_n(x_0)}{(x - x_0)^n - 0}
  \;=\; \frac{r_n'(\xi_1)}{n\,(\xi_1 - x_0)^{n-1}}.
\]
Applying the same theorem to $r_n'$ and $(\,\cdot\, - x_0)^{n-1}$ on
the interval between $x_0$ and~$\xi_1$, using $r_n'(x_0) = 0$, gives
$\xi_2$ strictly between $x_0$ and~$\xi_1$ with
\[
  \frac{r_n'(\xi_1)}{n\,(\xi_1 - x_0)^{n-1}}
  \;=\; \frac{r_n''(\xi_2)}{n(n-1)\,(\xi_2 - x_0)^{n-2}}.
\]
Continuing for $n - 1$ steps yields $\xi_{n-1}$ strictly between
$x_0$ and~$x$ with
\[
  \frac{r_n(x)}{(x - x_0)^n}
  \;=\; \frac{r_n^{(n-1)}(\xi_{n-1})}{n!\,(\xi_{n-1} - x_0)}.
\]
As $x \to x_0$, every $\xi_k \to x_0$, and by the definition of the
$n$-th derivative (applied to $r_n^{(n-1)}$),
\[
  \lim_{\xi \to x_0} \frac{r_n^{(n-1)}(\xi) - r_n^{(n-1)}(x_0)}{n!\,(\xi - x_0)}
  \;=\; \frac{r_n^{(n)}(x_0)}{n!} \;=\; 0,
\]
using $r_n^{(n-1)}(x_0) = 0$ and $r_n^{(n)}(x_0) = 0$.
Therefore $r_n(x)/(x - x_0)^n \to 0$, i.e.\ $r_n(x) = o((x - x_0)^n)$.
\end{proof}

\begin{proof}[\Cref{thm:taylor-lagrange}]
\proofanchor{taylor-lagrange}
Fix~$x \neq x_0$ and define on the interval between $x_0$ and~$x$:
\[
  g(t) \;=\; f(x) - \sum_{k=0}^{n} \frac{f^{(k)}(t)}{k!}\,(x - t)^k,
  \qquad
  h(t) \;=\; (x - t)^{n+1}.
\]
Both $g$ and $h$ are differentiable on the chosen interval; computing
$g'$ telescopically (every term except the last cancels) gives
\[
  g'(t) \;=\; -\frac{f^{(n+1)}(t)}{n!}\,(x - t)^n,
  \qquad
  h'(t) \;=\; -(n+1)(x - t)^n.
\]
At the endpoints $t = x_0$ and $t = x$,
\[
  g(x_0) = f(x) - \sum_{k=0}^{n} \frac{f^{(k)}(x_0)}{k!}\,(x - x_0)^k,
  \qquad
  g(x) = 0,
  \qquad
  h(x_0) = (x - x_0)^{n+1},
  \qquad
  h(x) = 0.
\]
Cauchy's generalized MVT applied to $(g, h)$ provides~$c$ strictly
between $x_0$ and~$x$ with
\[
  \frac{g(x_0) - g(x)}{h(x_0) - h(x)} \;=\; \frac{g'(c)}{h'(c)}
  \;=\; \frac{-f^{(n+1)}(c)\,(x - c)^n / n!}{-(n+1)(x - c)^n}
  \;=\; \frac{f^{(n+1)}(c)}{(n+1)!}.
\]
Equivalently
$g(x_0) = \frac{f^{(n+1)}(c)}{(n+1)!}\,(x - x_0)^{n+1}$, which
rearranges to the displayed Taylor formula.
\end{proof}

\begin{proof}[\Cref{cor:second-deriv-test}]
\proofanchor{second-deriv-test}
Apply Taylor with Peano remainder of order~$2$ at~$x_0$ using
$f'(x_0) = 0$:
\[
  f(x) - f(x_0) \;=\; \tfrac{1}{2} f''(x_0)\,(x - x_0)^2 + o\bigl((x - x_0)^2\bigr).
\]
If $f''(x_0) > 0$, the leading $\tfrac{1}{2}f''(x_0)(x - x_0)^2$ is
strictly positive for $x \neq x_0$ near~$x_0$ and dominates the
$o((x - x_0)^2)$ term, so $f(x) > f(x_0)$ on a punctured
neighbourhood — a strict local minimum. Symmetric reasoning gives
the maximum case. When $f''(x_0) = 0$ both signs are possible
($f(x) = x^4$ has a minimum, $f(x) = -x^4$ has a maximum, $f(x) = x^3$
has neither — all with vanishing first and second derivatives at~$0$).
\end{proof}

\subsection*{Examples}

\begin{example}[Extrema of a cubic]
\label{ex:cubic-extrema}
Let $f(x) = x^3 - 3x^2$. Then $f'(x) = 3x^2 - 6x = 3x(x - 2)$, so the
critical points (where $f'$ vanishes) are $x = 0$ and $x = 2$.
$f''(x) = 6x - 6$, giving $f''(0) = -6 < 0$ and $f''(2) = 6 > 0$.
By \cref{cor:second-deriv-test}, $f$ has a strict local maximum at
$x = 0$ (with $f(0) = 0$) and a strict local minimum at $x = 2$
(with $f(2) = -4$).
\end{example}

\begin{example}[Taylor expansion of $e^x$]
\label{ex:taylor-exp}
Take $f(x) = e^x$. All derivatives are $f^{(k)}(x) = e^x$, so
$f^{(k)}(0) = 1$. Taylor with Lagrange remainder
(\cref{thm:taylor-lagrange}) at $x_0 = 0$ gives, for every
$x \in \R$ and every~$n$,
\[
  e^x \;=\; \sum_{k=0}^{n} \frac{x^k}{k!} \;+\; \frac{e^c}{(n+1)!}\,x^{n+1}
\]
for some $c$ strictly between $0$ and~$x$. Bounding $e^c \le e^{|x|}$:
\[
  \Bigl| e^x - \sum_{k=0}^{n} \frac{x^k}{k!} \Bigr|
  \;\le\; \frac{e^{|x|} \,|x|^{n+1}}{(n+1)!}.
\]
The right-hand side tends to~$0$ as $n \to \infty$ for any
fixed~$x$, recovering the familiar series
$e^x = \sum_{k=0}^{\infty} x^k / k!$.
\end{example}

\begin{example}[Taylor for limits]
\label{ex:taylor-limit}
Compute $\lim_{x \to 0} (\sin x - x) / x^3$. The Peano-remainder
expansion $\sin x = x - x^3/6 + o(x^3)$ gives
\[
  \frac{\sin x - x}{x^3} \;=\; \frac{-x^3/6 + o(x^3)}{x^3}
  \;=\; -\tfrac{1}{6} + o(1) \;\to\; -\tfrac{1}{6}.
\]
The same limit can be obtained by three applications of l'Hôpital,
but Taylor expansion turns the calculation into one step.
\end{example}

\begin{problem}
\label{prob:taylor-estimate-e}
Use the Taylor polynomial of $e^x$ about $x_0 = 0$ to compute
$e \approx \sum_{k=0}^{n} 1/k!$ for some $n$, and find an $n$ that
guarantees four-decimal accuracy
($|e - \sum_{k=0}^{n} 1/k!| < 5 \times 10^{-5}$).
\end{problem}

\begin{proof}[Solution]
By \cref{ex:taylor-exp} at $x = 1$,
\[
  \Bigl| e - \sum_{k=0}^{n} \frac{1}{k!} \Bigr| \;\le\; \frac{e^1}{(n+1)!}
  \;\le\; \frac{3}{(n+1)!},
\]
using the crude bound $e \le 3$. We need $(n+1)! \ge 3 / (5 \times 10^{-5}) = 60000$.
Computing factorials: $7! = 5040$, $8! = 40320$, $9! = 362880$. So
$(n+1)! \ge 60000$ first holds at $n + 1 = 9$, i.e.\ $n = 8$.

The estimate is therefore
$e \approx \sum_{k=0}^{8} 1/k! = 1 + 1 + 1/2 + 1/6 + 1/24 + 1/120 + 1/720 + 1/5040 + 1/40320 \approx 2.71828$,
agreeing with $e = 2.71828\ldots$ to all five decimals shown. The
bound $3/9! \approx 8.3 \times 10^{-6}$ is comfortably below the
required tolerance.
\end{proof}
